% ─────────────────────────────────────────────────────────────────────────────
\section{Double flexion (\S{}6.1.6)}
% ─────────────────────────────────────────────────────────────────────────────

\begin{maitre}
Pour une panne de toiture inclin\'ee, dans quel plan la charge
s'exerce-t-elle~?
\end{maitre}

\begin{apprenti}
La charge est verticale (poids, neige), mais la panne est inclin\'ee.
Donc la charge cr\'ee \`a la fois un moment dans le plan de la section
(axe fort $y$) \textbf{et} hors du plan (axe faible $z$)~?
\end{apprenti}

\begin{maitre}
Exactement. C'est la double flexion.
La d\'ecomposition de la charge $q_d$ selon les axes de la panne est~:
\end{maitre}

\begin{equation}
q_y = q_d \cdot \cos(\alpha) \qquad \text{(axe fort)}
\qquad
q_z = q_d \cdot \sin(\alpha) \qquad \text{(axe faible)}
\label{eq:decomp}
\end{equation}

\noindent o\`u $\alpha$ est la pente de la toiture. Les moments correspondants~:

\begin{equation}
M_{y,d} = \frac{q_y \cdot L^2}{8}, \qquad M_{z,d} = \frac{q_z \cdot L^2}{8}
\end{equation}

La v\'erification EN~1995 \S{}6.1.6 pour la double flexion donne
deux crit\`eres (les deux doivent \^etre satisfaits)~:

\begin{align}
\frac{\sigma_{m,y,d}}{f_{m,d}} + k_m \cdot \frac{\sigma_{m,z,d}}{f_{m,d}} &\leq 1{,}0
\label{eq:biaxial1}\\
k_m \cdot \frac{\sigma_{m,y,d}}{f_{m,d}} + \frac{\sigma_{m,z,d}}{f_{m,d}} &\leq 1{,}0
\label{eq:biaxial2}
\end{align}

\noindent avec $k_m = 0{,}7$ pour les sections rectangulaires.

\begin{lstlisting}[style=python, caption={Décomposition de charge et interaction biaxiale}]
import math

pente_rad = math.radians(config.pente_deg)

# Décomposition de la charge de calcul selon les axes de la panne
q_y = q_d * math.cos(pente_rad)   # composante axe fort
q_z = q_d * math.sin(pente_rad)   # composante axe faible

# Moments correspondants (poutre simplement appuyée)
M_y = q_y * L**2 / 8.0   # kN.m
M_z = q_z * L**2 / 8.0   # kN.m

# Contraintes (MPa) — même conversion que §6.1.6
sigma_y = M_y * 1e3 / materiau.W_cm3   / 10.0
sigma_z = M_z * 1e3 / materiau.W_z_cm3 / 10.0

# Deux taux d'interaction — §6.1.6 (les deux doivent être <= 1.0)
k_m = 0.7   # sections rectangulaires
taux_1 = sigma_y / f_m_d + k_m * sigma_z / f_m_d
taux_2 = k_m * sigma_y / f_m_d + sigma_z / f_m_d
\end{lstlisting}

\begin{maitre}
\`A quoi sert ce $k_m = 0{,}7$~?
\end{maitre}

\begin{apprenti}
Il r\'eduit la contribution de l'axe faible dans le crit\`ere ?
\end{apprenti}

\begin{maitre}
Oui, et inversement. C'est un facteur de r\'eduction de l'interaction~:
la norme reconnait que les deux moments n'atteignent pas leur valeur maximale
au m\^eme point de la section en m\^eme temps.
Pour une section carr\'ee ($b = h$), $k_m$ serait \'egal \`a 1{,}0
et les deux crit\`eres seraient identiques.
\end{maitre}